\documentclass[11pt,a4paper]{article}

% --- Packages ---
\usepackage[utf8]{inputenc}
\usepackage[margin=0.75in]{geometry}
\usepackage{lmodern}
\usepackage{amsmath,amssymb,amsfonts}
\usepackage{graphicx}
\usepackage{xcolor}
\usepackage{enumitem}
\usepackage{titlesec}
\usepackage{hyperref}

% --- Color Definitions ---
\definecolor{navyblue}{RGB}{0,32,96}
\definecolor{darkgray}{RGB}{60,60,60}

% --- Hyperlink Setup ---
\hypersetup{
    colorlinks=true,
    linkcolor=navyblue,
    urlcolor=navyblue
}

% --- Section Formatting ---
\titleformat{\section}{\color{navyblue}\large\bfseries}{}{0em}{}[\color{navyblue}\titlerule]
\titlespacing*{\section}{0pt}{8pt}{4pt}

\setlist[itemize]{leftmargin=1.5em, itemsep=2pt, topsep=2pt}

% --- Document ---
\begin{document}

\pagestyle{empty} % Single-page document, suppress page numbers

% Header
\begin{center}
    {\LARGE \bfseries \color{navyblue} Minimax Theory \& Robust Clustering Estimators for Coherent Scene Understanding} \\[0.4em]
    {\large \textbf{Executive Summary / Pre-Proposal Inquiry}} \\[0.3em]
    \textbf{Program:} ONR Code 311 (Mathematical and Data Sciences) --- Automated Image Understanding \\
    \textbf{PI:} Dr. Soham Jana, Assistant Professor (\href{mailto:sjana@nd.edu}{sjana@nd.edu}) \\
    Department of Applied and Computational Mathematics and Statistics (ACMS), University of Notre Dame \\
    \href{https://janasoham.github.io/}{janasoham.github.io}
\end{center}

\vspace{-0.5em}

\section*{1. Motivation \& DoD Relevance}
Automated perception and decision-making systems operating in Naval environments face extreme noise, coherent speckle, partial occlusion, and adversarial disruption. Standard computer vision heuristics and deep network models lack formal finite-sample guarantees, leading to high error rates under heavy-tailed or dynamic noise. This project bridges mathematical statistics and visual perception by framing image understanding as a non-asymptotic estimation problem on high-dimensional graphs, establishing optimal recovery limits for complex scene structures.

\section*{2. Theoretical Innovations \& Methodology}
Our framework delivers mathematically rigorous, computationally efficient estimators for image segmentation, scene graph parsing, and multi-sensor data fusion:
\begin{itemize}
    \item \textbf{Non-Asymptotic Minimax Guarantees:} Establishing finite-sample minimax risk bounds for graph-based image estimation, providing provable limits on target recovery.
    \item \textbf{Adversarial \& Heavy-Tailed Robustness:} Formulating distance-based and spectral clustering estimators that maintain stability under corrupted observation graphs, missing edges, and heavy-tailed speckle noise.
    \item \textbf{Computational Scalability:} Translating high-dimensional matrix and graph estimation into fast $O(N \log N)$ algorithms suitable for real-time scene partitioning.
\end{itemize}

\section*{3. System Architecture \& Pipeline}
\begin{figure}[h!]
    \centering
    % Replace the box below with your graphic file: \includegraphics[width=0.92\textwidth]{pipeline.pdf}
    \framebox[\textwidth][c]{\rule{0pt}{1.2in} \color{darkgray}\small [ Visual Pipeline Diagram: Raw Noisy Stream $\rightarrow$ High-Dim Graph $\rightarrow$ Robust Estimator $\rightarrow$ Provable Recovery ]}
    \caption{\small Proposed estimation pipeline mapping raw spatial/spectral data to provably robust scene understandings.}
\end{figure}

\vspace{-1em}

\section*{4. Projected Deliverables \& Scope}
\begin{itemize}
    \item \textbf{Theoretical Foundations:} Unified minimax bounds for community detection and non-parametric estimation in dynamic visual graphs.
    \item \textbf{Algorithmic Suite:} Open-source, scalable algorithms for robust clustering under corrupted or incomplete image data.
    \item \textbf{Validation:} Benchmark evaluation on high-density spatial arrays and coherent imaging datasets.
\end{itemize}

\section*{5. Proposed Effort \& Budget Scale}
\begin{itemize}
    \item \textbf{Duration:} 3 Years (Target Start: Fall 2027)
    \item \textbf{Estimated Budget:} $\sim \$150,000$ -- $\$180,000$ / year (Supports 1 Ph.D. Student, PI summer effort, and conference travel).
\end{itemize}

\end{document}
